\section{Research Roadmap and Current Implementation Status}
\label{sec:roadmap}

This roadmap reflects the current state of the project after UGV01 bench
instrumentation, powered motion bring-up, and the complete smooth/rough benign
square-route baseline collection. The project is no longer blocked on hardware
bring-up or nominal data collection. The core software stack, embedded
telemetry path, timestamped logging pipeline, 20-run benign matrix,
reproducible run-level split, and real-log attack replay are
operational. The frozen alarm and uncertainty definitions and the complete
multi-start offline statistical attack campaign are now also operational. The
remaining work is to complete the matched strategic-attacker comparison,
evaluate transport transfer, add independent prospective validation, and
generate final capability maps and manuscript results.

The project is approximately 90\% complete overall: engineering implementation
is approximately 95\% complete, experimental analysis approximately 90\%, and
publication preparation approximately 80--85\%.

\subsection{Current Completed Foundation}

The following components are already implemented and tested in the current
codebase:

\begin{itemize}
    \item A nonlinear mobile digital-twin pipeline with tracked-drive-compatible
    kinematics, EKF prediction and update, innovation covariance computation,
    Mahalanobis/NIS detection, detectability-bound calculation, stealth-bound
    calculation, and confidence/envelope logic.

    \item Synthetic experiment support for step bias, coordinate freeze, value
    replay, slow drift, strategic scheduling, latency perturbation, CSV
    logging, plotting, ROC analysis, and empirical probability-of-detection
    summaries.

    \item UGV01 embedded telemetry through the active firmware path
    \texttt{ugv01\_gps\_dev}, with BN220 GPS integrated on the verified
    original RX wiring path.

    \item Combined UGV01 telemetry through \texttt{T:147}, including GPS, IMU,
    voltage, encoder counts, sequence number, firmware sample time, firmware
    send time, GPS validity, satellite count, HDOP, and NMEA health counters.

    \item Edge-side logging with source/sample time, send time, edge-arrival
    time, estimate time, alarm time placeholder, clock-offset calibration,
    packet-loss count, request-failure count, stale-packet flag, queue depth,
    and HTTP latency.

    \item Real-log replay utilities that pass hardware telemetry through the
    digital-twin EKF and detector path.

    \item Pre-battery analysis scripts for stationary GPS summaries, bench
    telemetry latency and health reports, hardware replay consistency review,
    and tracked-rover calibration preparation.

    \item Dataset-oriented square-route logging scripts with run metadata for
    route, surface, speed, network condition, trial number, attack type, square
    turn settings, and straight-command settings.
\end{itemize}

\subsection{Current Dataset Status}

The primary moving route for the benign baseline is a repeated
$0.5\,\mathrm{m}$ square loop executed three times per trial
(\texttt{square0p5x3}). The current surface definitions are:

\begin{itemize}
    \item \textbf{Smooth surface:} kitchen floor.
    \item \textbf{Rough surface:} permeable concrete.
\end{itemize}

The baseline Wi-Fi dataset is complete:

\begin{center}
\begin{tabular}{lccc}
\hline
\textbf{Condition} & \textbf{Speed} & \textbf{Valid Trials} & \textbf{Status} \\
\hline
Permeable concrete, baseline Wi-Fi & Low & 5/5 & Complete \\
Permeable concrete, baseline Wi-Fi & Medium & 5/5 & Complete \\
Kitchen floor, baseline Wi-Fi & Low & 5/5 & Complete \\
Kitchen floor, baseline Wi-Fi & Medium & 5/5 & Complete \\
\hline
\end{tabular}
\end{center}

Thus, the full baseline subset is complete with
\[
2 \;\text{speeds} \times 2 \;\text{surfaces} \times 1 \;\text{transport condition}
\times 5 \;\text{trials} = 20 \;\text{valid runs}.
\]

No additional live baseline run is required for the current offline replay
stage. Additional rover runs should be reserved for a larger held-out set,
route-transfer evaluation, or final deployment validation.

The buffered-delay condition should be generated as an edge-side transport
stress condition or replay transformation rather than by changing rover
firmware behavior. This preserves the primary A1 threat boundary and avoids
making rover motion less repeatable during live data collection.

\subsection{Phase 1: Complete Benign Baseline Collection}

\textbf{Goal.} Complete the nominal hardware dataset needed for threshold
locking, detector calibration, and later attack replay.

\textbf{Completed tasks.}

\begin{itemize}
    \item Collected five low-speed and five medium-speed runs on both the
    smooth kitchen floor and rough permeable concrete.

    \item Screened each run for completion, GPS validity, sequence continuity,
    stale packets, request failures, route plausibility, and whether operator
    intervention changed the run semantics.

    \item Archived invalid or interrupted runs and excluded them from threshold
    locking and held-out evaluation.

    \item Generated a checksum-backed benign-run manifest and candidate audit
    with surface, speed, route, transport condition, trial, source filename,
    validity fields, and exclusion evidence.
\end{itemize}

\textbf{Exit criteria.}

\[
2 \;\text{surfaces} \times 2 \;\text{speeds} \times
1 \;\text{baseline Wi-Fi condition} \times 5 \;\text{trials}
= 20 \;\text{valid live benign runs}.
\]

\subsection{Phase 2: Replay, Quality Screening, and Train/Validation/Test Splits}

\textbf{Goal.} Convert the live benign dataset into a reproducible analysis
artifact.

\textbf{Tasks.}

\begin{itemize}
    \item Replay each valid benign run through the digital-twin pipeline.

    \item Verify that EKF state, innovation, covariance, NIS/Mahalanobis score,
    confidence, envelope region, timing fields, packet-loss fields, and queue
    fields remain populated.

    \item Produce nominal residual and NIS calibration plots separately for
    smooth and rough surfaces and for low and medium speed.

    \item Freeze development, validation, and held-out evaluation splits by
    complete run rather than by individual row.
\end{itemize}

\textbf{Exit criteria.} A versioned benign dataset and split manifest exist,
and every accepted run can be regenerated from raw logs through replay scripts.

\textbf{Current status.} This exit criterion is met. The accepted manifest
contains 20 checksum-identified runs split by complete trial into 12
development, four validation, and four held-out test runs.

\subsection{Phase 3: Threshold Locking and Baseline Detector Suite}

\textbf{Goal.} Lock detector thresholds using benign data only before any
attack evaluation.

\textbf{Tasks.}

\begin{itemize}
    \item Lock the main NIS/Mahalanobis threshold at a matched run-level
    false-alarm target:
    \[
        P_{\mathrm{FA}} \leq 0.05.
    \]

    \item Report per-update NIS calibration separately from run-level
    operational false alarms.

    \item Implement or finalize baseline detectors:
    fixed GPS residual threshold, fixed digital-twin residual threshold,
    fixed-covariance EKF--NIS, innovation-gated EKF, robust measurement
    handling, sequential drift detector, naive adaptive EKF,
    GPS-independent adaptive EKF, and evidence-gated digital twin.

    \item Freeze threshold files, detector configurations, and split metadata
    before attack replay.
\end{itemize}

\textbf{Exit criteria.} All detector variants have frozen thresholds at matched
false-alarm rates, and no attack data have been used to tune those thresholds.

\textbf{Current status.} The operational alarm policy is frozen for the current
offline study. It uses robust five-fix GPS initialization, motion-gated mission
start, and a three-of-five persistent NIS rule. Because the earlier trial-5
results were inspected during alarm diagnosis, all 20 benign runs form the
design corpus and false alarms are estimated with complete-run leave-one-out
evaluation. The naturally anomalous run remains included and produces the
single leave-one-out alarm, giving the preregistered point estimate
$P_{\mathrm{FA}}=1/20=0.05$ for every variant. This meets the empirical target
but still requires a confidence interval and is not an independent prospective
validation claim.

\subsection{Phase 4: Covariance-Poisoning Study}

\textbf{Goal.} Test the central security claim that residual-coupled adaptive
uncertainty can become an attack surface under slow semantic GPS drift.

\textbf{Core paired replay design.}

Each physical benign log should be replayed under:

\begin{enumerate}
    \item clean GPS,
    \item drifted GPS with fixed covariance,
    \item drifted GPS with naive residual-coupled adaptive covariance,
    \item drifted GPS with frozen clean covariance,
    \item drifted GPS with GPS-independent covariance prediction,
    \item drifted GPS with evidence-gated adaptation.
\end{enumerate}

\textbf{Primary measurements.}

\begin{itemize}
    \item Change in $\mathrm{tr}(Q_k)$ and $\mathrm{tr}(S_k)$.
    \item Change in NIS/Mahalanobis score.
    \item Detection delay.
    \item Maximum undetected reported-state cross-track error.
    \item Time spent above the mission tolerance before alarm.
    \item Gate-pass rate and covariance-escalation rate.
\end{itemize}

\textbf{Exit criteria.} The project can show, with paired runs on the same
physical logs, whether naive adaptation inflates covariance under slow drift
and whether GPS-independent or evidence-gated adaptation reduces the
harmful-but-stealthy region.

\textbf{Current status.} This analysis is complete for the current design
corpus. Against the primary frozen-clean control, naive residual-coupled
adaptation increases the attacked/clean $Q$ ratio by $0.0243$ with a 95\%
physical-run-clustered interval of $[0.0185, 0.0300]$ and decreases the
attacked/clean NIS ratio by $0.0135$ with interval
$[-0.0174, -0.0090]$. The covariance-poisoning mechanism is therefore
measurable. However, the additional undetected state error is only
$0.013\,\mathrm{m}$, detection probability is unchanged, and the harmful-run
probability difference is not statistically established. The correct current
claim is that the mechanism exists, while operational attacker advantage is
not yet established. The evidence gate admits residual feedback during about
36.6\% of updates and is not itself activated more often by the tested drift.

\subsection{Phase 5: Directional Detectability and Attack Curves}

\textbf{Goal.} Estimate how attack detectability changes with direction,
surface, speed, and transport stress.

\textbf{Attack profiles.}

\begin{itemize}
    \item Step bias with magnitudes
    \[
        0.5,\;1,\;2,\;3,\;5,\;7.5,\;10 \;\mathrm{m}.
    \]
    \item Coordinate freeze.
    \item Value replay.
    \item Slow drift.
    \item Strategic drift under matched amplitude, rate, injection-count, and
    horizon budgets.
\end{itemize}

\textbf{Tasks.}

\begin{itemize}
    \item Estimate $\epsilon_{50}$, $\epsilon_{90}$, and $\epsilon_{95}$ for
    along-track, cross-track, and least-favorable attack directions.

    \item Compare empirical detection curves against the theoretical
    direction-dependent target-detection amplitude and worst-case bound:
    \[
        \epsilon_{\mathrm{worst}}(k)
        =
        \sqrt{\lambda^\star \lambda_{\max}(S_k)}.
    \]

    \item Report harmful-but-stealthy probability, maximum undetected
    cross-track error, detection delay, and time above
    $\epsilon_{\mathrm{req}}$ before alarm.

    \item Compare all detectors at matched run-level false-alarm rate.
\end{itemize}

\textbf{Exit criteria.} The paper has real-data attack curves showing where
attacks are trivially detected, harmful-but-stealthy, or ambiguous.

\textbf{Current status.} The offline attack campaign on the existing corpus is
complete. It evaluates 24 attack profiles at three post-motion injection times
for five detector variants across all 20 accepted physical runs, producing
7,200 paired attacked scenarios. Directional $\epsilon_{50}$, $\epsilon_{90}$,
and $\epsilon_{95}$ estimates, physical-run-clustered confidence intervals,
detection delays, and harmful-but-stealthy probabilities are generated. These
results describe the current design corpus; a separate prospective dataset is
still required for independent generalization claims.

\subsection{Phase 6: Strategic Attacker and Evidence-Gate Evaluation}

\textbf{Goal.} Evaluate whether an informed attacker gains additional stealth by
injecting drift during naturally high-uncertainty windows, and whether
evidence-gated adaptation limits that advantage.

\textbf{Tasks.}

\begin{itemize}
    \item Compare oblivious slow drift against white-box strategic drift under
    the same amplitude, rate, direction, injection-count, and horizon budgets.

    \item Keep any hindsight-oracle attacker separate as a non-realistic upper
    bound.

    \item Measure strategic advantage in terms of alarm delay, covariance
    escalation, gate-pass rate, and maximum undetected deviation.
\end{itemize}

\textbf{Exit criteria.} The manuscript can state whether evidence gating
meaningfully reduces strategic uncertainty harvesting.

\subsection{Phase 7: Transfer, Timing Stress, and Capability Maps}

\textbf{Goal.} Show how the locked detector configuration transfers across
surface, speed, route, and transport stress without retuning.

\textbf{Tasks.}

\begin{itemize}
    \item Apply the edge-side delay and jitter emulator to generate stressed
    transport conditions from accepted live logs.

    \item Report p50/p95/p99 and maximum edge-pipeline latency, source-to-alarm
    latency when clock-offset quality permits, packet loss, stale-packet rate,
    queue depth, CPU use, and memory footprint.

    \item Evaluate transfer on held-out routes if a reliable figure-eight or
    additional route-reference setup is available.

    \item Generate capability maps classifying operating conditions as
    Detectable, Ambiguous, or Low Sensitivity relative to the mission tolerance
    $\epsilon_{\mathrm{req}}$.
\end{itemize}

\textbf{Exit criteria.} The project has safety-anchored capability maps rather
than only detector-score plots.

\subsection{Phase 8: Artifact Freeze and Manuscript Assembly}

\textbf{Goal.} Freeze the complete experimental artifact package and assemble a
conference-ready manuscript.

\textbf{Tasks.}

\begin{itemize}
    \item Freeze raw logs, processed logs, attack labels, train/validation/test
    splits, threshold files, configuration files, scripts, lockfiles, and a
    checksum manifest.

    \item Regenerate every table and figure from the frozen artifact package.

    \item Assemble the manuscript around the central claim:
    residual-coupled covariance adaptation can become an attack surface;
    GPS-independent features block direct covariance poisoning; and
    evidence-gated adaptation limits strategic uncertainty harvesting.

    \item Include a claims-versus-evidence table that maps each paper claim to
    the figure, table, dataset, or theorem supporting it.

    \item Clearly separate synthetic validation, stationary hardware logs,
    moving benign hardware logs, replay-injected attacks, and any live attack
    demonstrations.
\end{itemize}

\textbf{Exit criteria.} A complete manuscript draft, frozen artifact package,
publication-quality figures, and final experimental logbook are available.

\subsection{Immediate Next Actions}

The next concrete steps are:

\begin{enumerate}
    \item Complete the matched ordinary-versus-strategic drift comparison under
    identical attack budgets.

    \item Expand the buffered delay/jitter condition and evaluate transfer without
    changing rover firmware or rerunning nominal motion.

    \item Collect a small untouched prospective benign validation set for the
    final frozen alarm and detector definitions.

    \item Generate final causal figures, capability maps, and the
    claims-versus-evidence table for manuscript assembly.
\end{enumerate}
